% Generated draft: requires mathematical and source review.
\documentclass[11pt]{article}
\usepackage[margin=25mm]{geometry}
\usepackage{fontspec}
\setmainfont{Latin Modern Roman}
\usepackage{amsmath,amssymb,mathtools,unicode-math}
\setmathfont{Latin Modern Math}
\usepackage{enumitem,needspace,xurl,hyperref}
\setlength{\parindent}{0pt}
\setlength{\parskip}{5pt}
\setlength{\emergencystretch}{3em}
\newcommand{\N}{\mathbb{N}}
\newcommand{\Z}{\mathbb{Z}}
\newcommand{\Q}{\mathbb{Q}}
\newcommand{\R}{\mathbb{R}}
\newcommand{\C}{\mathbb{C}}
\newcommand{\F}{\mathbb{F}}
\newcommand{\A}{\mathbb{A}}
\newcommand{\PP}{\mathbb{P}}
\newcommand{\E}{\mathbb{E}}
\newcommand{\eps}{\varepsilon}
\newcommand{\dd}{\,\mathrm{d}}
\begin{document}
\section*{UnsolvedMath ID: 30006598}
\textbf{Source problem code:} catalogue entry
\subsection{Definitions and mathematical statement}
Let \(C\) and \(D\) be irreducible components of Lusztig nilpotent varieties of Dynkin type \(A\), in the common representation-theoretic setting intended by the supplied source. Write \(\pi(C)\) and \(\pi(D)\) for their associated simple modules. For a module \(M\), its socle \(\operatorname{soc}(M)\) is the sum of its simple submodules.

The supplied conjecture asserts that, for every such pair \(C,D\),
\[
 C\text{ is rigid}\quad\Longrightarrow\quad
 \operatorname{soc}\bigl(\pi(C)\times\pi(D)\bigr)
 \cong \pi(C*D).
\]
Here the displayed isomorphism expresses the source's equality of representations up to isomorphism. In particular, the assertion says that this socle is simple and has the isomorphism class attached to \(C*D\). Rigidity is required only of \(C\); the supplied statement imposes no rigidity hypothesis on \(D\).

\textbf{Unresolved formulation qualifications.} The record does not define rigidity, the assignment \(C\mapsto\pi(C)\), the ordered component operation \(*\), or the representation product \(\times\). It also does not specify the underlying field, ambient module category, dimension-vector conventions, or compatibility conditions between the varieties containing \(C\) and \(D\). These symbols retain their intended source-specific meanings; their definitions and conventions must be recovered from the cited sources before this becomes a fully specified standalone conjecture. The order of both products is retained exactly as supplied.
\subsection{Short English statement}
The conjecture says that if the first Lusztig component \(C\) is rigid, then the product of the representations attached to \(C\) and \(D\) has a simple socle, namely the representation attached to the ordered component product \(C*D\). The precise representation-theoretic constructions are missing from the supplied record and require review.
\subsection{Sources}
\begin{itemize}
\item Supplied local UnsolvedMath registry record, permanent ID \(30006598\), title \emph{Socles of Products from Lusztig Components}, source problem code \texttt{catalogue entry}, in the collection \emph{Oberwolfach Reports Open Problems}. The record lists proposed year 2026 and supplies no named proposer.
\item \emph{Representation Theory of Quivers and Finite-Dimensional Algebras} (2026), Oberwolfach Reports, Report 6/2026, pp.~345--432, DOI 10.4171/owr/2026/6. Supplied report and literature-source link: \url{https://doi.org/10.4171/owr/2026/6}. The extracted passage says that the conjecture was made in reference [LM25]; the record does not supply a complete bibliographic entry under that label.
\item Supplied literature source: Lapid--Minguez, \emph{binary operation on nilpotent varieties II}, \url{https://arxiv.org/abs/2111.05162}.
\item Literature assessment dated exactly \texttt{2026-08-22}: the supplied curation labels the problem open, says that the socle identity is still called a conjecture in the 2026 report, and says that the primary paper verifies only special cases, including a stated rigidity class. This is supplied dated curation, not an independently verified literature assessment; the special cases and rigidity class are not specified in the record.
\item The registry labels the statement \texttt{corrected\_verified} and reports that it was reconstructed and checked against the official Oberwolfach report. This is supplied verification metadata, not an independent verification of this draft or of mathematical correctness.
\end{itemize}
\end{document}
